\section{Limitations and Threats to Validity}\label{sec:limits}
\answers{Bounds how far the answers to RQ1--RQ7 generalise}

\paragraph*{Version drift} Scores, probes and bench cells refer to the versions listed in \texttt{VERSIONS.md}, accessed on 15~October 2026. Harnesses release weekly, and several were renamed or discontinued during the study; the matrix should be re-run rather than cited as current.

\paragraph*{Models and compute} No vendor API keys or logins were used, so the bench runs a single small, local, open-weight model through documented custom-endpoint features and says nothing about any vendor's flagship model. The GPU is shared with the owner's other work; a planned second model family was cancelled to release it, so the goal of at least two models per harness is \emph{not met}, and the scenario set and runs per cell were reduced, leaving wide intervals. So small a model also limits how often any injection is acted upon at all. Claude Code was routed to a non-Claude model, which its vendor does not support (Section~\ref{sec:method}).

\paragraph*{Headless only} Probes and bench runs are headless. Interactive trust dialogs, IDE extensions and GUI-only products (the Antigravity IDE, for instance) were not tested. A harness whose interactive mode is safer is under-credited, which is the conservative error for CI use.

\paragraph*{Benign-canary ceiling} The payloads are inert and plainly worded; disguised or adaptive attacks may succeed where these fail, so the bench rates are lower bounds.

\paragraph*{Black-box limits} Without TLS inspection we cannot say what an allowed host receives, and login-gated harnesses were observed only before login.

\paragraph*{Rating and classification} Both matrix raters worked from one shared fact base, so $\kappa$ measures the application of the rubric rather than the reliability of the evidence, and it excludes the N/A-versus-score cells (all \nDisagreeEOne{} disagreements, those included, are published). Incident classes and root causes had a single LLM rater and no agreement measurement, and RC1 is the residual class of an ordered procedure. Several scores turn on rubric interpretation, notably of classifier-gated defaults.

\paragraph*{Selection} The incident survey admits only publicly disclosed, harness-level advisories, most of them with CVE or GHSA identifiers; model-level failures are rarely tracked in this way, and agent-assisted search favours recent, well-indexed advisories. The RQ3 shares therefore describe disclosed harness defects, not overall risk.

\paragraph*{Corpora} E6 is synthetic and labelled by the author; E2's code-search sample favours popular, recently indexed repositories; E5 has one session per condition.
